\hypertarget{buxe1o-cuxe1o-thu-houx1ea1ch-ux111ux1ecbnh-hux1b0ux1edbng-nghux1ec1-nghiux1ec7p-kux1ef9-nux103ng-vuxe0-tux1b0-duy-thux1ef1c-tux1ebf-trong-doanh-nghiux1ec7p}{%
\section{Báo cáo Thu hoạch: Định hướng Nghề nghiệp, Kỹ năng và Tư duy
Thực tế trong Doanh
nghiệp}}

\hypertarget{mux1ee5c-tiuxeau-sux1ef1-kiux1ec7n}{%
\subsection{1. Mục tiêu Sự
kiện}}

\begin{itemize}
\tightlist
\item
  Giúp sinh viên năm 3 - 4 gạt bỏ những lầm tưởng về môi trường công sở,
  từ đó trang bị một hệ tư duy thực tế và nhạy bén nhằm thích ứng với
  những đòi hỏi khắt khe của thị trường lao động hiện đại.
\item
  Làm rõ bản chất thực chiến của các vị trí kỹ thuật như DevOps, Kỹ sư
  Dữ liệu (Data Engineer), đồng thời hiểu rõ các chuẩn mực vận hành
  chuyên nghiệp tại những tập đoàn đa quốc gia lớn.
\end{itemize}

\hypertarget{diux1ec5n-giux1ea3-khuxe1ch-mux1eddi}{%
\subsection{2. Diễn giả Khách
mời}}

\begin{itemize}
\tightlist
\item
  \textbf{Anh Trong H. Truong} - DevOps Engineer @ Endava Vietnam.
\item
  \textbf{Anh Danh Hoàng Hiếu Nghị} - AI Engineer @ AWS Community
  Builder \& AWS Student Builder Group Leader.
\item
  \textbf{Anh Đạt Phạm} - Data Analytics Engineer @ Kamereo \&
  Colgate-Palmolive.
\item
  \textbf{Anh Cường Nguyễn} - Process Engineer @ Colgate-Palmolive.
\end{itemize}

\hypertarget{nhux1eefng-ux111iux1ec3m-nhux1ea5n-quan-trux1ecdng-key-highlights}{%
\subsection{3. Những Điểm Nhấn Quan Trọng (Key
Highlights)}}

\hypertarget{thux1ef1c-tux1ebf-khux1ed1c-liux1ec7t-vuxe0-bux1ea3n-chux1ea5t-cuxf4ng-viux1ec7c-devops-trong-doanh-nghiux1ec7p}{%
\subsubsection{Thực tế khốc liệt và bản chất công việc DevOps trong
doanh
nghiệp}}

Thực tế công việc DevOps hoàn toàn không hào nhoáng như những dòng mô tả
trên lý thuyết. Đó là một chuỗi ngày phải chịu trách nhiệm trực on-call
luân phiên, xử lý sự cố hệ thống lúc nửa đêm, cho đến việc giải quyết
những bài toán hóc búa về tối ưu hóa chi phí đám mây và phân định rõ
ràng quyền sở hữu hệ thống.

Một bài học sâu sắc rút ra là công cụ (tools) chỉ mang tính thời điểm.
Giá trị cốt lõi bền vững chính là nền tảng tư duy hệ thống và khả năng
tự động hóa những tác vụ lặp đi lặp lại nhằm giải phóng sức lao động cho
toàn đội ngũ.

\hypertarget{sux1ef1-linh-houx1ea1t-cux1ee7a-kux1ef9-sux1b0-dux1eef-liux1ec7u-theo-tux1eebng-lux129nh-vux1ef1c-kinh-doanh}{%
\subsubsection{Sự linh hoạt của Kỹ sư Dữ liệu theo từng lĩnh vực kinh
doanh}}

Vai trò của một Data Analytics Engineer không hề cố định mà biến đổi vô
cùng linh hoạt theo đặc thù doanh nghiệp. Tại Kamereo (mô hình F\&B),
trọng tâm công việc là tối ưu chuỗi cung ứng hằng ngày, thiết kế các
dashboard trực quan để ban giám đốc đưa ra quyết định kinh doanh (GMV)
tức thời.

Trái lại, ở một tập đoàn sản xuất FMCG toàn cầu như Colgate-Palmolive,
người kỹ sư phải trực tiếp xuống nhà máy làm việc với hệ thống dữ liệu
IoT của máy móc, thiết bị sản xuất. Mục tiêu là để nâng cao hiệu suất
vận hành và thúc đẩy tiến trình chuyển đổi số dài hạn của nhà máy.

\hypertarget{buxe0i-hux1ecdc-ux111uxfac-kux1ebft-key-takeaways}{%
\subsection{4. Bài Học Đúc Kết (Key
Takeaways)}}

\hypertarget{tux1b0-duy-phuxe1t-triux1ec3n-sux1ef1-nghiux1ec7p-qua-5-giai-ux111oux1ea1n}{%
\subsubsection{Tư duy phát triển sự nghiệp qua 5 giai
đoạn}}

Tôi thực sự bị thuyết phục bởi tư duy ``thăng tiến năng lực thay vì
thăng tiến chức danh''. Bản thân mỗi người trẻ cần định vị rõ mình đang
ở đâu trên tháp năng lực: từ một \textbf{Follower} (chỉ biết làm theo
checklist), tiến lên thành \textbf{Learner} (chủ động đặt câu hỏi sâu
sắc để tích lũy thực chiến), bứt phá thành \textbf{Problem Solver} (cam
kết chất lượng đầu ra), trở thành \textbf{System Thinker} (nhìn nhận cục
diện để tối ưu lâu dài), và cuối cùng là \textbf{Super Star} (người dẫn
dắt và định hướng tầm nhìn cho cả tổ chức).

\hypertarget{nux1ec1n-tux1ea3ng-kux1ef9-thuux1eadt-vuxe0-huxe0nh-truxecnh-8-bux1b0ux1edbc-tux1ef1-thuxe2n}{%
\subsubsection{Nền tảng kỹ thuật và hành trình 8 bước tự
thân}}

Đối với mảng kỹ thuật chuyên sâu như DevOps, bài học lớn nhất là phải
kiên trì rèn luyện các kiến thức nền tảng như Linux, mạng máy tính, ngôn
ngữ lập trình (Python/Golang), Git/CI-CD và Container, thay vì chạy đua
theo các công cụ liên tục thay đổi.

Hành trình 8 bước truyền cảm hứng của anh Danh Hoàng Hiếu Nghị --- từ
một cậu sinh viên tò mò dấn thân vào chương trình First Cloud Journey,
miệt mài thực hành các bài Lab và làm đồ án, xây dựng Portfolio chất
lượng để trở thành đối tác AWS, và cuối cùng là tinh thần ``Share-Back''
để nâng đỡ thế hệ kế cận --- đã chứng minh rằng: sự bền bỉ và tinh thần
phụng sự cộng đồng chính là chìa khóa mở ra những cơ hội lớn.

\hypertarget{vux103n-huxf3a-doanh-nghiux1ec7p-vuxe0-triux1ebft-luxfd-ux111uxfang-viux1ec7c}{%
\subsubsection{Văn hóa doanh nghiệp và triết lý ``Đúng
Việc''}}

Tôi vô cùng ấn tượng với văn hóa ``No-Blame Post-Mortem'' tại các tập
đoàn công nghệ. Ở đó, khi xảy ra sự cố nghiêm trọng, mọi người cùng ngồi
lại phân tích lỗ hổng hệ thống để cải tiến, tuyệt đối không dùng để đổ
lỗi cho cá nhân.

Triết lý ``Wakon Yosai'' (Hòa hồn Dương tài) của người Nhật khuyên chúng
ta giữ gìn bản sắc cốt lõi nhưng phải chủ động làm chủ các kỹ nghệ tiên
tiến toàn cầu. Song song đó, triết lý ``Đúng Việc'' hướng mỗi người đến
việc gánh vác trách nhiệm thông qua ba cột trụ vững chắc: Làm Người
(quản trị nội tâm), Làm Nghề (phụng sự bằng thực lực) và Làm Dân (trách
nhiệm xã hội).

\hypertarget{ux1ee9ng-dux1ee5ng-vuxe0o-thux1ef1c-tux1ebf-applying-to-work}{%
\subsection{5. Ứng Dụng Vào Thực Tế (Applying to
Work)}}

\begin{itemize}
\tightlist
\item
  \textbf{Ngừng sao chép vô thức}: Bỏ ngay thói quen copy các dòng lệnh
  vô hồn trên mạng. Thay vào đó, phải luôn đào sâu bản chất và đặt câu
  hỏi ``Tại sao?'' trước khi tìm hiểu ``Làm như thế nào?''.
\item
  \textbf{Kỹ năng Data Storytelling}: Tập trung rèn luyện khả năng
  truyền thông dữ liệu, biến các bảng biểu số liệu khô khan thành những
  câu chuyện có cấu trúc để thúc đẩy hành động thực tế.
\item
  \textbf{Làm chủ vấn đề}: Chủ động nhận diện ``ai là chủ nhân thực sự
  của vấn đề'' và cam kết giải quyết triệt để, thay vì chỉ làm việc đối
  phó hoặc thụ động chờ đợi hướng dẫn.
\item
  \textbf{Tinh thần tập thể}: Thực hành văn hóa hỗ trợ đồng đội, đơn
  giản hóa các quy trình vận hành và tài liệu hóa (document) mọi giải
  pháp để chia sẻ tri thức cho toàn đội.
\end{itemize}

\hypertarget{trux1ea3i-nghiux1ec7m-khuxf3a-hux1ecdc-event-experience}{%
\subsection{6. Trải Nghiệm Khóa Học (Event
Experience)}}

\begin{itemize}
\tightlist
\item
  \textbf{Bài học từ sự chân thành}: Những câu chuyện thật, những trải
  nghiệm thực chiến đầy gai góc của các diễn giả đã giúp tôi dập tắt
  những ảo tưởng màu hồng của lý thuyết sách vở.
\item
  \textbf{Góc nhìn công nghệ thực tế}: Bản đồ công cụ DevOps đồ sộ ban
  đầu làm tôi choáng ngợp, nhưng cũng chỉ ra một lối đi rõ ràng: phải
  tập trung làm chủ nền tảng và hiểu cặn kẽ cách một ứng dụng thực sự
  vận hành trên môi trường production.
\item
  \textbf{Sử dụng công cụ thông minh}: Ý thức được việc sử dụng AI như
  một người trợ lý đắc lực để nâng cao năng suất cá nhân, nhưng tuyệt
  đối không được để AI làm thui chột đi tư duy chủ động và khả năng phản
  biện của chính mình.
\item
  \textbf{Kết nối cộng đồng}: Không khí thảo luận cởi mở và nhiệt huyết
  tại buổi meetup đã thắp lên trong tôi khát khao học hỏi, dấn thân xây
  dựng cộng đồng và liên tục nâng cấp tiêu chuẩn của bản thân.
\end{itemize}

\hypertarget{tux1ed5ng-kux1ebft-lessons-learned}{%
\subsection{7. Tổng Kết (Lessons
Learned)}}

\begin{itemize}
\tightlist
\item
  Việc mù quáng sao chép dòng lệnh mà không hiểu bản chất hệ thống chỉ
  tạo ra những lỗ hổng bảo mật và sự cố nguy hiểm trong tương lai.
\item
  Một kỹ sư DevOps xuất sắc không phải là một ``người hùng cô đơn'' cứu
  thế giới lúc nửa đêm, mà là người biết giao tiếp hiệu quả, phối hợp
  với đội ngũ và xây dựng một kiến trúc ổn định lâu dài.
\item
  Năng lực kỹ thuật chỉ giúp ta bước chân vào nghề, nhưng chính tư duy
  hệ thống và các kỹ năng mềm (giao tiếp, phản biện) mới là đôi cánh đưa
  ta tiến xa trên lộ trình sự nghiệp.
\end{itemize}

\emph{(Hình ảnh sự kiện có thể được bổ sung tại đây)}
